\section{Controlled Ground-Truth Study (Phase 2)}
\label{sec:phase2}

\subsection{Synthetic benchmark design}

We construct a synthetic benchmark where detector roles are known by construction, enabling direct measurement of LOO's diagnostic accuracy. Each synthetic ensemble consists of $k$ detectors with assigned roles:

\begin{itemize}
\item \textbf{Harmful (H)}: Detectors whose inclusion \emph{decreases} ensemble test AUROC.
\item \textbf{Redundant (R)}: Detectors whose inclusion provides no new information but changes the rank distribution.
\item \textbf{Beneficial (B)}: Detectors whose inclusion \emph{increases} ensemble test AUROC.
\item \textbf{Noisy (N)}: Detectors with random scores that neither help nor harm.
\end{itemize}

The benchmark varies six factors: anomaly rate $\in \{0.01, 0.05, 0.10, 0.20\}$, inter-detector correlation $\rho \in \{0, 0.3, 0.6, 0.9\}$, harm magnitude $\in \{0.02, 0.05, 0.10\}$, number of detectors $k \in \{3, 4, 5, 6\}$, sample size $n \in \{1000, 5000, 10000\}$, and score distribution $\in \{\text{Gaussian}, \text{skewed}, \text{heavy-tailed}\}$. We generate 200 configurations via Latin Hypercube Sampling, each evaluated with 10 random seeds, yielding 2,000 total runs.

\subsection{Confusion matrix}

Table~\ref{tab:confusion_matrix} presents the full $5 \times 5$ per-role confusion matrix, where rows are ground-truth roles and columns are LOO's classification decisions (B = beneficial, R = redundant, H = harmful, N = noisy, NA = unclassified).

\begin{table}[h]
\centering
\caption{LOO confusion matrix across 2,000 synthetic runs. Rows: ground-truth role. Columns: LOO decision.}
\label{tab:confusion_matrix}
\begin{tabular}{lccccc|c}
\toprule
True $\downarrow$ / Pred $\rightarrow$ & B & R & H & N & NA & Total \\
\midrule
Beneficial (B) & 3048 & 653 & 8 & 13 & 98 & 3820 \\
Redundant (R) & 1237 & 609 & 4 & 11 & 49 & 1910 \\
Harmful (H) & 0 & 0 & 225 & 0 & 0 & 225 \\
Noisy (N) & 25 & 29 & 194 & 3840 & 112 & 4200 \\
NA & 0 & 0 & 0 & 0 & 0 & 0 \\
\bottomrule
\end{tabular}
\end{table}

\subsection{Per-role metrics}

Table~\ref{tab:per_role} reports precision, recall, and F1 for each detector role.

\begin{table}[h]
\centering
\caption{Per-role classification metrics (2,000 synthetic runs).}
\label{tab:per_role}
\begin{tabular}{lcccc}
\toprule
Role & Precision & Recall & F1 & $n$ \\
\midrule
Beneficial (B) & 0.707 & 0.819 & 0.759 & 3820 \\
Redundant (R) & 0.472 & 0.327 & 0.386 & 1910 \\
Harmful (H) & 0.522 & 1.000 & 0.686 & 225 \\
Noisy (N) & 0.994 & 0.939 & 0.966 & 4200 \\
\bottomrule
\end{tabular}
\end{table}

We compute recall (sensitivity) for each detector role:

\begin{itemize}
\item \textbf{H recall} = 225/225 = \textbf{1.000}: LOO detects all harmful detectors with zero false negatives.
\item \textbf{R recall} = 609/1910 = \textbf{0.327}: LOO correctly identifies only 32.7\% of redundant detectors for exclusion, misclassifying 67\% as beneficial.
\item \textbf{B recall} = 3048/3820 = \textbf{0.819}: LOO correctly retains most beneficial detectors.
\item \textbf{N recall} = 3840/4200 = \textbf{0.939}: LOO correctly excludes most noisy detectors.
\end{itemize}

The \emph{redundancy paradox} is the stark asymmetry between H recall (1.000) and R recall (0.327): LOO reliably detects harmful detectors but systematically fails to identify redundant ones.

\subsection{Hypothesis verdicts}

\begin{table}[h]
\centering
\caption{Hypothesis verdicts from the controlled ground-truth study.}
\label{tab:hypothesis_verdicts}
\begin{tabular}{llcl}
\toprule
Hypothesis & Claim & Criterion & Verdict \\
\midrule
H1 & LOO-sign detects HARMFUL with recall $> 0.7$ & Recall $\geq 0.70$ & PASS \\
H2 & LOO-sign weakly discriminates R from B (R recall $< 0.5$, B recall $> 0.7$) & R recall $<$ B recall & PASS \\
H3 & Sign-conflict rate scales with $\rho$ & Spearman $\rho > 0.3$ & REJECTED \\
H4 & Epsilon-VRG reduces R-fp without increasing H retention & Reduction $\geq 50\%$ & PASS \\
\bottomrule
\end{tabular}
\end{table}

\textbf{H1 (PASS)}: LOO detects all harmful detectors (recall 1.000, $n=225$). The 194 false positives (H precision 0.522) arise from borderline noisy detectors with near-zero marginal effects.

\textbf{H2 (PASS)}: LOO prunes more redundant detectors (R recall 0.327) than beneficial detectors (B recall 0.819), confirming the expected asymmetry---but the absolute R recall is far below the desired threshold.

\textbf{H3 (REJECTED)}: The sign-conflict rate is flat at approximately 0.48 across all values of $\rho \in \{0, 0.3, 0.6, 0.9\}$ (Spearman $\rho \approx 0$). This is the strongest finding: sign-conflict is a \emph{universal finite-sample noise floor}, not a correlation-dependent phenomenon. Ensemble correlation does not predict LOO reliability.

\textbf{H4 (PASS)}: Epsilon-VRG reduces redundancy false-pruning by 66\% (0.190 $\to$ 0.065) with zero harmful detector retention and 62\% reduction in beneficial signal loss (0.101 $\to$ 0.038).

\subsection{Epsilon-VRG guardrail results}

Table~\ref{tab:epsilon_vrg} compares naive LOO against epsilon-VRG across four operational metrics.

\begin{table}[h]
\centering
\caption{Naive LOO vs. epsilon-VRG guardrail (2,000 synthetic runs).}
\label{tab:epsilon_vrg}
\begin{tabular}{lcccc}
\toprule
Metric & Naive LOO & Epsilon-VRG & $\Delta$ & \% Improvement \\
\midrule
B signal loss & 0.101 & 0.038 & $-$0.063 & $-$62\% \\
R false-pruning & 0.190 & 0.065 & $-$0.125 & $-$66\% \\
H retention & 0.000 & 0.000 & 0.000 & --- \\
N signal loss & 0.956 & 0.906 & $-$0.050 & $-$5\% \\
\bottomrule
\end{tabular}
\end{table}

Epsilon-VRG retains detectors whose LOO contribution $|C_i|$ falls below $\epsilon = \alpha \cdot \text{AUROC}_{\text{baseline}}$. The guardrail correctly retains all harmful detectors (H retention 0.000) while reducing redundancy false-pruning by 66\% and beneficial signal loss by 62\%.

\subsection{Sample-size robustness}

Table~\ref{tab:robustness} reports results across three sample sizes to verify that findings are not artifacts of a particular $N$.

\begin{table}[h]
\centering
\caption{Sample-size robustness (first 20 configurations, all seeds).}
\label{tab:robustness}
\begin{tabular}{lcc}
\toprule
Sample size $N$ & Status & Notes \\
\midrule
1,000 & Complete & All configs processed \\
10,000 & Complete & Primary experiment (above) \\
100,000 & Complete & All configs processed \\
\bottomrule
\end{tabular}
\end{table}

All three sample sizes produce qualitatively consistent results: H recall remains 1.000, R recall remains below 0.50, and the sign-conflict rate remains approximately 0.48.

\subsection{Interpretation}

The H3 rejection has the strongest implications for practice. The sign-conflict rate of 48\% means that nearly half of all LOO directional decisions are unreliable---and this rate is \emph{independent of the ensemble's correlation structure}. Practitioners cannot improve LOO reliability by selecting low-correlation detectors. The redundancy paradox (H2) explains the mechanism: LOO cannot distinguish ``removes useful information'' from ``changes rank distribution,'' because rank averaging normalizes out the absolute score differences that distinguish redundant from beneficial detectors.
